\documentclass[11pt,a4paper]{article}

\usepackage{amsmath}

\begin{document}
\tableofcontents

\section{Mécanique 1 - 1989}

\subsection{Chapitre 9 - Interaction deux particules}

\subsubsection{Calcul 1}

\begin{align*}
\frac{\partial r}{\partial x_1} &= \frac{\partial}{\partial r}([(x_2-x_1)^2+(y_2-y_1)^2+(z_2-z_1)^2]^\frac{1}{2})\\
&=\frac{\partial}{\partial r}((x_2-x_1)^2+(y_2-y_1)^2+(z_2-z_1)^2)\times-\frac{1}{2[(x_2-x_1)^2+(y_2-y_1)^2+(z_2-z_1)^2]}\\
&=-\frac{x_2-x_1}{r}
\end{align*}

\subsubsection{Notes}
D'après le PFD et la troisième loi de Newton\\
$
\frac{d^2(m_1r_1+m_2r_2)}{dt^2}=0\\
$
On pose
$
r_G=OG=\frac{m_1r_1+m_2r_2}{m_1+m_2}
$

\end{document}